\documentclass[11pt]{article}
\usepackage[margin=0.86in]{geometry}
\usepackage[T1]{fontenc}
\usepackage[utf8]{inputenc}
\usepackage{amsmath,amssymb,amsthm}
\usepackage{graphicx}
\usepackage{xcolor}
\usepackage{microtype}
\usepackage[colorlinks=true,urlcolor=blue!55!black,citecolor=blue!55!black]{hyperref}

\newtheorem{theorem}{Theorem}
\newtheorem{corollary}[theorem]{Corollary}
\newcommand{\R}{\mathbb R}
\newcommand{\D}{\mathcal X}

\title{An interval-length boundary for the open spectral-set\\
number-rigidity question}
\author{Candidate substantial partial answer to the question of arXiv:1702.07323}
\date{11 August 2026}

\begin{document}
\maketitle

\begin{center}
\fbox{\parbox{0.92\textwidth}{\textbf{Status: candidate partial result, likely valid.}
For a widely separated countable union of intervals, we give an exact
interval-length criterion for \emph{linear} number rigidity.  Exponentially
decaying interval lengths produce genuinely number-rigid open spectral sets
with infinitely many components.  Polynomially decaying lengths produce
explicit open finite-measure sets whose DPPs are not linearly number rigid.
Whether the latter are nonlinearly number rigid remains open.}}
\end{center}

\section*{Source question}

On PDF page 3 of \emph{Rigid stationary determinantal processes in
non-Archimedean fields}, Qiu asks \cite{source}:
\begin{quote}
Let $S\subset\R$ be an open subset with finite positive Lebesgue measure.
Is the determinantal point process on $\R$ induced by the correlation kernel
$\widehat{\mathbf 1_S}(x-y)$ number rigid?
\end{quote}
The complete source page is reproduced below.

\begin{center}
\includegraphics[width=0.94\textwidth]{source_question_page.png}
\end{center}

The question is universal: it asks about every such open set.  Our result
settles a broad positive subfamily and identifies an explicit family for which
all $L^2$ linear-statistic methods necessarily fail.

\section*{The structure factor}

For a measurable $S\subset\R$ of finite positive measure, let $\D_S$ be the
stationary projection DPP with kernel
\[
 K_S(x,y)=\widehat{\mathbf 1_S}(x-y).
\]
Write $\rho=|S|$.  Up to Fourier-normalization constants and the harmless
rescaling to unit intensity, the continuous covariance spectral density, or
structure factor, is
\begin{equation}\label{eq:structure}
 s_S(t)=\rho-|S\cap(S+t)|=\frac12|S\mathbin\triangle(S+t)|.
\end{equation}
Indeed, Plancherel and the convolution theorem give
$\widehat{|K_S(0,\cdot)|^2}(t)=|S\cap(S+t)|$.  Multiplying $s_S$ by a positive
constant does not affect any reciprocal-integrability statement.

We use the precise conclusion that is safe in the later criterion of
Lachi\`eze-Rey \cite{linear}: in dimension one, a stationary DPP is
\emph{linearly} number rigid on bounded intervals exactly when
\begin{equation}\label{eq:criterion}
 \int_{0}^{\varepsilon}\frac{dt}{s_S(t)}=\infty.
\end{equation}
See PDF page 19, especially the sentence following Corollary 2.  The paper
explicitly cautions on PDF page 8 that the word ``linear'' is sometimes
omitted but all its methods concern linear rigidity.  Divergence in
\eqref{eq:criterion} is nevertheless enough for ordinary Ghosh--Peres number
rigidity, since $L^2$ linear reconstruction is an actual measurable
reconstruction.  Convergence proves only failure of linear number rigidity.

\section*{Exact classification for separated intervals}

Let $(\ell_n)_{n\geq1}$ satisfy
\[
 0<\ell_n\leq1,\qquad \sum_{n\geq1}\ell_n<\infty,
\]
and define the open set
\begin{equation}\label{eq:S}
 S_{\ell}=\bigcup_{n\geq1}(3n,3n+\ell_n).
\end{equation}

\begin{theorem}\label{thm:classification}
For $0<|t|<1$,
\begin{equation}\label{eq:omega}
 s_{S_{\ell}}(t)=\omega_{\ell}(|t|),\qquad
 \omega_{\ell}(u):=\sum_{n\geq1}\min\{u,\ell_n\}.
\end{equation}
Consequently $\D_{S_{\ell}}$ is linearly number rigid if and only if
\begin{equation}\label{eq:lengthcriterion}
 \int_0^\varepsilon\frac{du}{\sum_n\min\{u,\ell_n\}}=\infty.
\end{equation}
Whenever \eqref{eq:lengthcriterion} holds, $\D_{S_{\ell}}$ is also number
rigid in the ordinary nonlinear sense.
\end{theorem}

\begin{proof}
For $|t|<1$, an interval in \eqref{eq:S} can meet only its own translate;
distinct three-unit blocks remain disjoint.  The overlap contributed by the
$n$th interval is $(\ell_n-|t|)_+$.  Thus
\[
 |S_{\ell}|-|S_{\ell}\cap(S_{\ell}+t)|
 =\sum_n\bigl(\ell_n-(\ell_n-|t|)_+\bigr)
 =\sum_n\min\{|t|,\ell_n\},
\]
which proves \eqref{eq:omega}.  The equivalence is \eqref{eq:criterion}.
If the integral diverges, linear reconstruction gives number rigidity for
bounded intervals.  For an arbitrary bounded Borel set $B$, choose a bounded
interval $I\supset B$: the outside of $B$ reveals $\D(I\setminus B)$, while
the outside of $I$ determines $\D(I)$, so
$\D(B)=\D(I)-\D(I\setminus B)$ is also determined.
\end{proof}

This formula is also the exact competition between the number of intervals
longer than the translation scale and the total mass of shorter intervals:
\[
 \omega_\ell(u)=u\,\#\{n:\ell_n\geq u\}
 +\sum_{\ell_n<u}\ell_n.
\]

\section*{Two opposite infinite-component regimes}

\begin{corollary}[A genuinely rigid open set]\label{cor:exp}
For $\ell_n=2^{-n}$, the open set $S_\ell$ has finite positive measure,
infinitely many components, and $\D_{S_\ell}$ is number rigid.
\end{corollary}

\begin{proof}
If $2^{-(N+1)}<u\leq2^{-N}$, then
\[
 \omega_\ell(u)=Nu+\sum_{n>N}2^{-n}=Nu+2^{-N}
 \asymp u\log(1/u).
\]
Hence \eqref{eq:lengthcriterion} diverges like
$\int_0 du/(u\log(1/u))$.  Theorem~\ref{thm:classification} applies.
\end{proof}

This supplies a particularly elementary positive infinite-union subcase.
Lin, Qiu, and Wang independently constructed more intricate rigid generalized
Cantor sets and complementary infinite unions \cite{cantor}; their PDF page 1
states the prior finite-union boundary and the goal of finding special
infinite-union examples.

\begin{corollary}[A sharp linear obstruction family]\label{cor:poly}
Fix $p>1$ and put $\ell_n=n^{-p}$.  Then $S_\ell$ is open and has finite
positive measure, but $\D_{S_\ell}$ is not linearly number rigid.  More
precisely,
\[
 s_{S_\ell}(t)\asymp |t|^{1-1/p}\qquad(t\to0).
\]
In particular, for $p=2$ one obtains the explicit set
\[
 S=\bigcup_{n\geq1}(3n,3n+n^{-2}),
 \qquad s_S(t)\asymp\sqrt{|t|}.
\]
\end{corollary}

\begin{proof}
Take $N=\lfloor u^{-1/p}\rfloor$.  Then
\[
 \omega_\ell(u)=Nu+\sum_{n>N}n^{-p}.
\]
Both terms are comparable to $u^{1-1/p}$, by the integral bounds for the
$p$-series tail.  Since $1-1/p<1$, the reciprocal is locally integrable.
The equivalence in Theorem~\ref{thm:classification} proves failure of linear
number rigidity.
\end{proof}

\section*{Deep upgrade attempt and exact obstruction}

The disjoint frequency intervals give an orthogonal decomposition of the
Fourier projection.  A tempting upgrade is therefore to regard
$\D_{S_\ell}$ as an unlabelled superposition of independent sine processes
and study mark erasure.  This route fails at its first probabilistic step:
orthogonal sum of projection \emph{operators} does not make the corresponding
spatial DPP the independent union of the component DPPs.  Expanding the
determinant of the summed kernel leaves mixed interference terms; independence
occurs only in special spatially block-diagonal situations.

A direct Palm/deletion-tolerance route also stops at a real obstruction.  The
reduced Palm kernel is a rank-one modification of the full projection, but
neither the interval-length formula nor the second-order structure factor
controls absolute continuity of its point-process law.  Thus convergence of
\eqref{eq:lengthcriterion} cannot honestly be promoted to nonlinear
non-rigidity.  This is not a cosmetic gap: arXiv:2409.18519 explicitly
separates linear from nonlinear rigidity, and known perturbed lattices show
that nonlinear rigidity can survive failure of linear rigidity.

Accordingly, Corollary~\ref{cor:poly} is a strong method barrier, not a full
counterexample to Qiu's question.  Settling it would require a new theorem
showing deletion tolerance, equivalence of different-order Palm laws, or an
independent nonlinear reconstruction mechanism.

\section*{Novelty and verification bounds}

Searches through 11 August 2026 covered the exact question, arbitrary open
spectral sets, infinite unions of intervals, structure-factor criteria,
superpositions of sine processes, Palm equivalence, and deletion tolerance.
No located source states Theorem~\ref{thm:classification} or the polynomial
obstruction family.  The decisive analytic criterion is from
arXiv:2409.18519, so novelty is claimed only for the interval-length reduction,
its phase boundary, and the explicit application to the source question.

The accompanying script checks the exact finite-truncation overlap formula
and the predicted exponential and polynomial asymptotics.  It is a sanity
check, not part of the proof.

\begin{thebibliography}{9}
\bibitem{source}
Y. Qiu, \emph{Rigid stationary determinantal processes in non-Archimedean
fields}, arXiv:1702.07323, 2017.  Question on PDF p.~3.

\bibitem{linear}
R. Lachi\`eze-Rey, \emph{Rigidity of random stationary measures and
applications to point processes}, arXiv:2409.18519v3, 2025.  Linear-rigidity
caveat on PDF p.~8; stationary-DPP criterion and linear conclusion on PDF
p.~19.

\bibitem{cantor}
Z. Lin, Y. Qiu, and K. Wang, \emph{Number rigid determinantal point processes
induced by generalized Cantor sets}, arXiv:2407.14168, 2024.
\end{thebibliography}

\end{document}
